% GENERATED FILE. Edit canonical lessons, then run npm run generate:books.
% canonical-source-hash: fnv1a64:cab1e5f5e1925be8
% canonical-lessons: GE-C422-kreuzfahrt, GE-C422-anlegen, GE-C422-ablegen, GE-C422-schaukeln, GE-C422-kapitaen

\chapter{Auf See}
\label{ch:ge-auf-see}
\hlchaptermodality{422}

\begin{chapteropening}
\textbf{By the end of this chapter:} \emph{I can talk about a cruise: the cruise, docking, casting off, rocking and the captain.}

Complete the last lesson of chapter 422: I can talk about a cruise: the cruise, docking, casting off, rocking and the captain.
\end{chapteropening}

\section[die Kreuzfahrt]{die Kreuzfahrt — the cruise}

\label{lesson:GE-C422-kreuzfahrt}



\begin{quote}
Before the new one: say the German for the life jacket, then the German for seasick.
\end{quote}

\subsection*{You'll want to know: die Kreuzfahrt}
\textbf{die Kreuzfahrt} — \textquotedblleft{}the cruise\textquotedblright{}.

\begin{quote}
Meine Großeltern machen im Sommer eine Kreuzfahrt.
\end{quote}

\begin{quote}
\emph{My grandparents are going on a cruise in the summer.}
\end{quote}

Plural: \emph{die Kreuzfahrten}. \emph{Das Kreuzfahrtschiff} — the cruise ship.

\begin{cousinweb}
\emph{Kreuzfahrt} is built from \emph{kreuzen} 'to cruise, to sail back and forth', from \emph{Kreuz} 'cross', which goes back to Latin \emph{crux} (Lewis \& Short), and \emph{Fahrt} 'voyage', from \emph{fahren} (EtymDB).
\end{cousinweb}

\begin{grammarlens}[title={What you've built}]
One more word for telling, explaining and describing at B1.
\end{grammarlens}

\subsection*{Your turn}
\begin{itemize}
\raggedright
  \item \emph{Say it:} \emph{die Kreuzfahrt}
  \item \emph{Say it:} \emph{die Kreuzfahrt}, once more
  \item \emph{Say it:} a sentence of your own with \emph{die Kreuzfahrt}
  \item \emph{Recall:} say the German for the life jacket, then the German for seasick, then say \emph{die Kreuzfahrt} again
\end{itemize}

\begin{tcolorbox}[breakable,colback=teal!4,colframe=teal!35!black,title={Before you move on}]
What does die Kreuzfahrt mean? (\textquotedblleft{}The cruise\textquotedblright{}.) Say it once more.
\end{tcolorbox}
\section[anlegen]{anlegen — to dock, to moor; to put on; to invest}

\label{lesson:GE-C422-anlegen}



\begin{quote}
Before the new one: say the German for seasick, then the German for the cruise.
\end{quote}

\subsection*{You'll want to know: anlegen}
\textbf{anlegen} — \textquotedblleft{}to dock, to moor; to put on; to invest\textquotedblright{}.

\begin{quote}
Das Schiff legt am Morgen im Hafen an.
\end{quote}

\begin{quote}
\emph{The ship docks in the harbour in the morning.}
\end{quote}

Separable: \emph{das Schiff legt an}, \emph{es hat angelegt}. \emph{Geld anlegen} — to invest money.

\begin{cousinweb}
\emph{anlegen} is built from \emph{an} 'on, at' and \emph{legen} 'to lay', inherited from Old High German \emph{leggen} (EtymDB): a boat laid alongside the quay.
\end{cousinweb}

\begin{grammarlens}[title={What you've built}]
One more word for telling, explaining and describing at B1.
\end{grammarlens}

\subsection*{Your turn}
\begin{itemize}
\raggedright
  \item \emph{Say it:} \emph{anlegen}
  \item \emph{Say it:} \emph{anlegen}, once more
  \item \emph{Say it:} a sentence of your own with \emph{anlegen}
  \item \emph{Recall:} say the German for seasick, then the German for the cruise, then say \emph{anlegen} again
\end{itemize}

\begin{tcolorbox}[breakable,colback=teal!4,colframe=teal!35!black,title={Before you move on}]
What does anlegen mean? (\textquotedblleft{}To dock, to moor; to put on; to invest\textquotedblright{}.) Say it once more.
\end{tcolorbox}
\section[ablegen]{ablegen — to cast off, to leave port; to take off (a coat); to take (an exam)}

\label{lesson:GE-C422-ablegen}



\begin{quote}
Before the new one: say the German for the cruise, then the German for to dock.
\end{quote}

\subsection*{You'll want to know: ablegen}
\textbf{ablegen} — \textquotedblleft{}to cast off, to leave port; to take off (a coat); to take (an exam)\textquotedblright{}.

\begin{quote}
Um zehn Uhr legt das Schiff wieder ab.
\end{quote}

\begin{quote}
\emph{At ten o'clock the ship casts off again.}
\end{quote}

Separable. \emph{Eine Prüfung ablegen} — to take an exam. \emph{Den Mantel ablegen} — to take off one's coat.

\begin{cousinweb}
\emph{ablegen} is built from \emph{ab} 'off, away' and \emph{legen} 'to lay', inherited from Old High German \emph{leggen} (EtymDB).
\end{cousinweb}

\begin{grammarlens}[title={What you've built}]
One more word for telling, explaining and describing at B1.
\end{grammarlens}

\subsection*{Your turn}
\begin{itemize}
\raggedright
  \item \emph{Say it:} \emph{ablegen}
  \item \emph{Say it:} \emph{ablegen}, once more
  \item \emph{Say it:} a sentence of your own with \emph{ablegen}
  \item \emph{Recall:} say the German for the cruise, then the German for to dock, then say \emph{ablegen} again
\end{itemize}

\begin{tcolorbox}[breakable,colback=teal!4,colframe=teal!35!black,title={Before you move on}]
What does ablegen mean? (\textquotedblleft{}To cast off, to leave port; to take off (a coat); to take (an exam)\textquotedblright{}.) Say it once more.
\end{tcolorbox}
\section[schaukeln]{schaukeln — to rock, to swing}

\label{lesson:GE-C422-schaukeln}



\begin{quote}
Before the new one: say the German for to dock, then the German for to cast off.
\end{quote}

\subsection*{You'll want to know: schaukeln}
\textbf{schaukeln} — \textquotedblleft{}to rock, to swing\textquotedblright{}.

\begin{quote}
Bei Wind schaukelt das Schiff hin und her.
\end{quote}

\begin{quote}
\emph{In the wind the ship rocks to and fro.}
\end{quote}

\emph{Die Schaukel} — the swing in a playground.

\begin{cousinweb}
\emph{schaukeln} is built on \emph{Schaukel} 'swing'. EtymDB lists the verb without a source, and the older history of the word is not settled.
\end{cousinweb}

\begin{grammarlens}[title={What you've built}]
One more word for telling, explaining and describing at B1.
\end{grammarlens}

\subsection*{Your turn}
\begin{itemize}
\raggedright
  \item \emph{Say it:} \emph{schaukeln}
  \item \emph{Say it:} \emph{schaukeln}, once more
  \item \emph{Say it:} a sentence of your own with \emph{schaukeln}
  \item \emph{Recall:} say the German for to dock, then the German for to cast off, then say \emph{schaukeln} again
\end{itemize}

\begin{tcolorbox}[breakable,colback=teal!4,colframe=teal!35!black,title={Before you move on}]
What does schaukeln mean? (\textquotedblleft{}To rock, to swing\textquotedblright{}.) Say it once more.
\end{tcolorbox}
\section[der Kapitän]{der Kapitän — the captain}

\label{lesson:GE-C422-kapitaen}



\begin{quote}
Before the new one: say the German for to cast off, then the German for to rock.
\end{quote}

\subsection*{You'll want to know: der Kapitän}
\textbf{der Kapitän} — \textquotedblleft{}the captain\textquotedblright{}.

\begin{quote}
Der Kapitän begrüßt die Gäste an Bord.
\end{quote}

\begin{quote}
\emph{The captain welcomes the passengers on board.}
\end{quote}

Plural: \emph{die Kapitäne}. Also the captain of a team.

\begin{cousinweb}
\emph{Kapitän} came into German through Middle High German \emph{kapitān} from Medieval Latin \emph{capitāneus} 'chief' (EtymDB), built on Latin \emph{caput} 'head' (Lewis \& Short).
\end{cousinweb}

\begin{grammarlens}[title={What you've built}]
One more word for telling, explaining and describing at B1.
\end{grammarlens}

\subsection*{Your turn}
\begin{itemize}
\raggedright
  \item \emph{Say it:} \emph{der Kapitän}
  \item \emph{Say it:} \emph{der Kapitän}, once more
  \item \emph{Say it:} a sentence of your own with \emph{der Kapitän}
  \item \emph{Recall:} say the German for to cast off, then the German for to rock, then say \emph{der Kapitän} again
\end{itemize}

\begin{tcolorbox}[breakable,colback=teal!4,colframe=teal!35!black,title={Before you move on}]
What does der Kapitän mean? (\textquotedblleft{}The captain\textquotedblright{}.) Say it once more.
\end{tcolorbox}
